% ============================================================================
%  Reconstruction-Quality-Targeted Sensor Placement for 1-D Salinity Fields
%  Catherine Sun -- MECH 493 final report
%  Format: JMLR-style single-column journal article
% ============================================================================
\documentclass[11pt,letterpaper]{article}

% ---------- Packages --------------------------------------------------------
\usepackage[margin=1.25in]{geometry}
\usepackage{amsmath,amssymb,amsfonts}
\usepackage{booktabs}
\usepackage{graphicx}
\usepackage{caption}
\usepackage{subcaption}
\usepackage{siunitx}
\usepackage{float}
\usepackage{enumitem}
\usepackage{xcolor}
\usepackage[round,authoryear]{natbib}
\usepackage[hidelinks,colorlinks=true,linkcolor=blue!50!black,
            citecolor=blue!50!black,urlcolor=blue!50!black]{hyperref}
\usepackage{url}
\usepackage{microtype}
\usepackage{titlesec}
\usepackage{array}

% ---------- Formatting (JMLR-style) -----------------------------------------
\titleformat{\section}{\large\bfseries}{\thesection.}{0.6em}{}
\titleformat{\subsection}{\normalsize\bfseries}{\thesubsection}{0.6em}{}
\titleformat{\subsubsection}{\normalsize\itshape}{\thesubsubsection}{0.5em}{}
\setlength{\parskip}{4pt}
\setlength{\parindent}{14pt}
\bibliographystyle{plainnat}

% ---------- Title block -----------------------------------------------------
\title{\bfseries Targeted Sensor Placement
for Salinity Fields}

\author{
  Catherine Sun \\
  \small University of British Columbia \\
}
\date{April 2026}

% ============================================================================
\begin{document}
\maketitle

% ---------- Abstract --------------------------------------------------------
\begin{abstract}
\noindent
Monitoring salinity along estuaries requires optimization of sensor placement and greedy
acquisition-function designs are the standard approach. However
, their optimality guaranties apply to local information-theoretic
objectives rather than to the downstream reconstruction error within
the scientifically critical sub-region. We introduce the \emph{Dataset
Quality Score} (DQS), a scale-free $[0,1]$ metric that compares any
sensor design against an arc-length-equidistant geometric benchmark
on a gradient-weighted RMSE in the steep region.

We consider a logarithmic test function calibrated to the Sooke
estuary and perform reconstruction using a Gaussian process with a
Mat\'ern-3/2 kernel. Four acquisition functions and three illustrative
baselines are evaluated along with a global differential-evolution optimizer
, on sensor counts $N \in \{3, \dots, 22\}$ with
observation noise $\sigma = 0.9$~PSU. Results are averaged over 30
seeds for the greedy methods and 10 for the optimiser.

The noise-averaged optimiser attains DQS $= 0.937$ at $N = 20$, while
the strongest greedy heuristic (MaxDist) reaches only $0.675$ --- a
structural gap of $0.26$ DQS, or roughly $28\%$ of achievable quality.
The gap peaks at $0.47$ in the mid-budget regime ($N = 10$) and
narrows as MaxDist gradually fills the halocline. The noiseless
optimiser saturates at DQS $= 1$ from $N = 6$, bounding how few
sensors suffice when placed optimally.

The framework generalises to any field whose scientific value
concentrates in an identifiable sub-region and quantifies the cost
of greedy one-step-ahead selection against a jointly optimised
reference.
\end{abstract}

\vspace{4pt}

\newpage
% ============================================================================
\section{Introduction}
\label{sec:intro}

\subsection{Motivation}
As a vital resource, water underpins the survival of organisms and the integrity of habitats in any given system. However, in Canada, the frequency of events that compromise the availability and quality of water is increasing. Climate change, through rising sea levels, altered nutrient distribution, and increasing water temperatures presents unprecedented challenges for water management. Pollution sources encourage harmful algae blooms and disrupt aquatic ecosystems through chemical contamination \cite{ASRESU2025109194, oneil}. Meanwhile, forest fires can impact drainage systems by introducing contaminants while floods wash sediments, nutrients, and pollutants into bodies of water \cite{bladon}. These threats highlight the growing need for monitoring and understanding water quality, both to safeguard ecosystems and respond to environmental changes. Continuous
monitoring of the salinity is essential
early warning detection, but is constrained by sensor cost, maintenance, and
data-link bandwidth. Operators typically deploy between three and twenty
fixed-point probes along a channel that spans several kilometres. The
practical question is therefore: \emph{where should the next sensor
go?}

\subsection{Problem statement}
Given a one-dimensional domain $\mathcal{X}\subset\mathbb{R}$ and a
target sensor count $N$, the task is to choose locations
$S=\{x_{1},\dots,x_{N}\}\subset\mathcal{X}$ that maximise the quality
of a field $\hat f_{S}$ reconstructed from noisy measurements
$y_{i}=f(x_{i})+\varepsilon_{i}$. Because reconstruction quality is
application-specific, the choice of scoring metric is itself a design
decision: in estuarine monitoring the scientifically relevant quantity
is the accuracy of $\hat f_{S}$ inside the steep region, while
errors on the flat offshore plateau carry less weight.

\subsection{Contributions}
This paper makes four contributions.
\begin{enumerate}[itemsep=2pt,topsep=2pt]
  \item We define the \emph{Dataset Quality Score} (DQS), a
        $[0,1]$-valued metric that normalises a problem-specific
        gradient-weighted RMSE against an arc-length-equidistant
        geometric benchmark.
  \item We position DQS within the targeted-IMSE tradition of
        \citet{williams2000sequential}, \citet{picheny2010adaptive},
        and \citet{bect2012sequential}, and within the
        regret-vs-best-known convention of \citet{frazier2018tutorial}.
  \item We empirically quantify the gap between greedy and globally
        optimised designs on a 1-D salinity test function calibrated to
        the Sooke estuary, separating noiseless ceilings from
        noise-averaged realities.
  \item We report calibration diagnostics (NLPD, $1\sigma$ coverage)
        as orthogonal secondary scores, demonstrating that the DQS
        rankings are not an artefact of miscalibrated posteriors.
\end{enumerate}

\subsection{Roadmap}
Section~\ref{sec:related} reviews the relevant literature on
experimental design, Gaussian processes, and acquisition functions.
Section~\ref{sec:problem} formalises the test function, domain, and
notation. Section~\ref{sec:dqs-framework} introduces the DQS framework.
Section~\ref{sec:setup} describes the evaluation protocols and
computational cost. Section~\ref{sec:results} reports the empirical
gap between greedy and optimised designs. Section~\ref{sec:calibration}
reports calibration diagnostics. Sections~\ref{sec:discussion}
and~\ref{sec:conclusions} discuss the implications and conclude.

% ============================================================================
\section{Background and Related Work}
\label{sec:related}

\subsection{The sensor-placement problem}
Coastal and estuarine salinity fields exhibit steep gradients near
freshwater inflows that relax into weakly varying offshore
plateaus~\citep{geyer2014estuarine}. Because monitoring budgets cap the
number of deployable sensors well below the spatial resolution of the
underlying physics, operators face a discrete design problem: choose
$N$ locations $S\subset\mathcal{X}$ to maximise the quality of the
reconstructed field. The problem is combinatorial
($\binom{|\mathcal{X}|}{N}$ candidate configurations) and, in its exact
form, NP-hard~\citep{ko1995exact}.

\subsection{Bayesian experimental design and information theory}
Classical optimal experimental
design~\citep{fedorov1972theory,chaloner1995bayesian} frames sensor
placement as maximisation of an information criterion. The A-, D-, and
E-optimal criteria target the expected posterior covariance directly;
the more flexible \emph{mutual information} (MI) criterion
\citep{caselton1984optimal,krause2008near} targets uncertainty
reduction between measured and unmeasured locations,
\begin{equation}
\label{eq:mi}
I(S;\bar S)\;=\;H(\bar S)-H(\bar S\mid S),
\end{equation}
with $\bar S=\mathcal{X}\setminus S$ and $H$ differential entropy.

\citet{krause2008near} proved that MI is submodular and monotone under
a mild regularity condition, so the greedy algorithm achieves a
$(1-1/e)$ approximation of the optimal MI:
\begin{equation}
\label{eq:greedy-bound}
I(S_{\text{greedy}})\;\ge\;\bigl(1-\tfrac{1}{e}\bigr)\,I(S^{*}).
\end{equation}
Equation~\eqref{eq:greedy-bound} underpins most acquisition-function
methods. Crucially, however, the guarantee concerns the
information-theoretic objective rather than a downstream reconstruction
metric, and it leaves open a gap of up to $\approx 37\%$ that a global
search might close.

\subsection{Gaussian-process surrogates}
The Gaussian process (GP)~\citep{rasmussen2006gaussian} plays two roles
here: (i) the reconstruction operator used to score any sensor set,
and (ii) the generative model whose posterior variance defines the
classical acquisition functions. A GP prior $f\sim\mathcal{GP}(\mu,k)$
with observations
$\mathbf{y}=\mathbf{f}+\boldsymbol{\varepsilon}$,
$\boldsymbol{\varepsilon}\sim\mathcal{N}(\mathbf{0},\sigma_n^{2}\mathbf{I})$,
yields the posterior
\begin{align}
\mu_{S}(x) &= \mu(x)+\mathbf{k}_{xS}
   \bigl[\mathbf{K}_{SS}+\sigma_n^{2}\mathbf{I}\bigr]^{-1}
   (\mathbf{y}-\boldsymbol{\mu}_{S}), \label{eq:gp-mean}\\
\sigma_{S}^{2}(x) &= k(x,x)-\mathbf{k}_{xS}
   \bigl[\mathbf{K}_{SS}+\sigma_n^{2}\mathbf{I}\bigr]^{-1}
   \mathbf{k}_{Sx}. \label{eq:gp-var}
\end{align}
The Mat\'ern-3/2 kernel
\begin{equation}
\label{eq:matern}
k(x,x')=\sigma_{f}^{2}\!\left(1+\tfrac{\sqrt{3}r}{\ell}\right)
\exp\!\left(-\tfrac{\sqrt{3}r}{\ell}\right),\qquad r=\|x-x'\|,
\end{equation}
is widely used for environmental data because it is once mean-square
differentiable, matching the mild roughness of geophysical fields
without imposing the smoothness of the squared-exponential
kernel~\citep{stein1999interpolation}.

Figure~\ref{fig:kernel} illustrates the effective correlation
$K(x_{\text{ref}},x)$ used in this work. Modelling salinity in
$\ln x$ rather than $x$ (right panel) preserves a near-stationary
correlation structure across the domain, so a single global length
scale captures both the steep near-mouth gradient and the offshore
plateau. \citet{herrera2022optimal} reports the same observation in
the air-pollution context: a log-transform on the spatial coordinate
turns a strongly non-stationary field into one whose correlation can
be modelled with a stationary kernel.

\begin{figure}[H]
\centering
\includegraphics[width=0.95\linewidth]{dqs_results/paper_fig_kernel.png}
\caption{Mat\'ern-3/2 correlation $K(x_{\text{ref}},x)$ for four
reference locations. Left: kernel evaluated in linear $x$. Right:
kernel evaluated in $\ln x$ (the formulation used in this work). The
log-space formulation matches the geometry of the test function
$f(x)=a\ln x+c$, allowing a single length scale to serve the steep
near-mouth region (red band) and the flat offshore plateau.}
\label{fig:kernel}
\end{figure}

\subsection{Acquisition functions}
Greedy sequential designs select each next sensor by optimising an
acquisition functional $a(\cdot;S_{t})$. The literature offers many
variants:
\begin{itemize}[itemsep=2pt,topsep=2pt]
  \item \textbf{MaxVar / UCB}~\citep{srinivas2010gaussian}: select
        $\arg\max_{x}\sigma_{S}(x)$.
  \item \textbf{Expected and Probability of
        Improvement}~\citep{mockus1978application,jones1998efficient}:
        select the point most likely to improve the best-so-far
        observation.
  \item \textbf{IVR and weighted IVR}
        \citep{sacks1989design,gramacy2009adaptive}: minimise a
        (weighted) integral of posterior variance reduction.
  \item \textbf{Mutual Information (GMI)}~\citep{krause2008near}: a
        closed-form surrogate for~\eqref{eq:mi}.
  \item \textbf{MaxDist}~\citep{johnson1990minimax}: maximin distance
        to the existing set; a geometry-only heuristic.
  \item \textbf{GradMaxVar}: a gradient-weighted posterior-variance
        rule that biases selection toward steep-gradient regions.
\end{itemize}
All of these rules optimise a \emph{local} surrogate; none directly
targets the reconstruction quality in the steep region.

\subsection{Reconstruction-quality metrics and benchmarks}
\label{sec:metrics-litreview}
The metric used to evaluate a sensor design is itself a methodological
choice with a well-developed literature. Three families dominate.

\subsubsection{Uniform integrated squared error}
The classical \emph{integrated mean squared prediction error} (IMSPE)
of \citet{sacks1989design} integrates the GP posterior variance
$\sigma_{S}^{2}(x)$ uniformly over the domain;
\citet{krause2008near} report root-mean-squared error (RMSE) under the
same uniform weighting. These metrics are appropriate when the
scientific value of the field is spatially homogeneous.

\subsubsection{Targeted (weighted) integrated squared error}
When interest concentrates on a sub-region or a level set,
\citet{williams2000sequential}, \citet{picheny2010adaptive}, and
\citet{bect2012sequential} replace the uniform integral with a weighted
one,
\begin{equation}
\label{eq:tIMSE}
\mathrm{tIMSE}(S)\;=\;\int_{\mathcal{X}}
\sigma_{S}^{2}(x)\,w(x)\,\mathrm{d}x,
\end{equation}
where $w(x)\ge 0$ encodes the region of interest. Picheny's $w(x)$ is a
tube indicator around a target level; Bect's is the probability of being
near a failure threshold. The gradient-weighted RMSE
$\mathcal{E}_{\text{steep}}$ defined in equation~\eqref{eq:steep-wrmse}
sits squarely in this targeted-IMSE tradition, with
$w(x)\propto|f'(x)|$ emphasising the steep region.

\subsubsection{Calibration metrics}
The negative log predictive density (NLPD) and continuous ranked
probability score (CRPS) jointly score mean accuracy and uncertainty
calibration~\citep{gneiting2007strictly,marchant2012bayesian,binois2019replication}.
NLPD is a strictly proper scoring rule but conflates calibration with
accuracy: a method with wide, well-calibrated predictive bands can
outscore one with tight bands close to truth. In monitoring
applications the latter is usually preferred, so NLPD is most commonly
reported as a secondary diagnostic alongside an error metric, not as
the primary design objective.

\subsubsection{Oracle- and regret-normalised scores}
Bayesian-optimisation benchmarking commonly reports performance as
\emph{simple regret} relative to the global optimum or to a best-known
reference~\citep{frazier2018tutorial}. The Dataset Quality Score
(Section~\ref{sec:dqs}) follows this normalisation tradition: the
arc-length-equidistant configuration of equation~\eqref{eq:arclength}
serves as a problem-specific reference, while differential evolution at
large $N$ provides an empirical upper bound that the clip at
$\mathrm{DQS}=1$ absorbs when the benchmark is itself sub-optimal.
Reporting both raw $\mathcal{E}_{\text{steep}}$ in PSU and the
normalised DQS preserves direct comparison with studies that report
only RMSE while adding the interpretability of a $[0,1]$ score.

\subsection{Global optimisation as an alternative}
When the acquisition objective is replaced by the reconstruction score
itself, the design problem becomes a continuous-variable black-box
optimisation over $\mathbb{R}^{N}$. Differential evolution
\citep{storn1997differential} is a population-based gradient-free
optimiser that has been applied successfully to non-convex placement
problems in hydrology and atmospheric
science~\citep{reed2013evolutionary}. Unlike greedy acquisition, DE can
revisit and redistribute previously committed sensors, allowing it to
approach the true optimum at the cost of
$\mathcal{O}(10^{4}\text{--}10^{5})$ objective evaluations.

\subsection{Gap addressed by this work}
Three gaps motivate this study: (i) \emph{metric alignment}---greedy
acquisition rules are tuned for variance or information, not for the
steep error; (ii) \emph{missing benchmark}---results are usually
reported as raw RMSE without a problem-specific upper bound; and
(iii) \emph{noise treatment}---noiseless and noise-averaged performance
are routinely conflated. The DQS framework introduced in
Section~\ref{sec:dqs-framework} closes all three gaps.

% ============================================================================
\section{Problem Formulation and Notation}
\label{sec:problem}

\subsection{Test function and domain}
\label{sec:test-function}
Salinity is modelled along a one-dimensional along-channel coordinate
$x\in[1,1400]\,\mathrm{m}$ by
\begin{equation}
\label{eq:testfun}
f(x)=a\ln(x)+c,\qquad a=1.49,\;\; c=13.62,
\end{equation}
fitted to transformed observations from the Sooke estuary dataset
(Appendix~\ref{app:data}). The logarithmic form reproduces the
near-mouth steep region and the flat offshore plateau with a single
closed-form expression and has analytic gradient $f'(x)=a/x$. The
domain is partitioned into a \emph{steep region}
$\mathcal{X}_{\text{steep}}=[1,200]\,\mathrm{m}$ and a \emph{flat region}
$\mathcal{X}_{\text{flat}}=[200,1400]\,\mathrm{m}$. A uniform evaluation
grid $\mathcal{G}=\{x_{j}\}_{j=1}^{M}$ with $M=1400$ is used for all
numerical integrals.

\subsection{Notation}
\label{sec:notation}
Symbols used throughout the paper are summarised in
Table~\ref{tab:notation}.

\begin{table}[H]
\centering
\caption{Notation.}
\label{tab:notation}
\small
\begin{tabular}{@{}ll@{}}
\toprule
Symbol & Meaning \\
\midrule
$\mathcal{X}=[1,1400]$\,m & 1-D along-channel domain \\
$\mathcal{X}_{\text{steep}}=[1,200]$\,m & Steep sub-region \\
$\mathcal{X}_{\text{flat}}=[200,1400]$\,m & Flat offshore sub-region \\
$\mathcal{G}=\{x_{j}\}_{j=1}^{1400}$ & Uniform evaluation grid \\
$f(x)=a\ln x + c$ & Test function ($a{=}1.49$, $c{=}13.62$) \\
$N$ & Number of sensors \\
$S=\{x_{1},\dots,x_{N}\}$ & A sensor design \\
$y_{i}=f(x_{i})+\varepsilon_{i}$ & Noisy observation \\
$\sigma$ & Observation noise std (set to 0.9 PSU) \\
$\hat f_{S}(x)$ & GP posterior mean given $S$ \\
$\sigma^{2}_{S}(x)$ & GP posterior variance given $S$ \\
$w_{j}\propto 1/x_{j}$ & Discrete gradient weight on $\mathcal{X}_{\text{steep}}$ \\
$\mathcal{E}_{\text{steep}}(S)$ & Gradient-weighted RMSE on $\mathcal{X}_{\text{steep}}$ \\
$S^{\star}$ & Arc-length-equidistant benchmark design \\
$\mathcal{E}_{\text{steep}}^{\text{ref}}=0.5594$ PSU
                                   & Benchmark reference error \\
$\mathrm{DQS}(S)$ & Dataset Quality Score, $\in[0,1]$ \\
$S_{\mathrm{DE}}(N)$ & DE-optimised design at size $N$ \\
\bottomrule
\end{tabular}
\end{table}

% ============================================================================
\section{The Dataset Quality Score Framework}
\label{sec:dqs-framework}

\subsection{Gradient-weighted RMSE on the region of interest}
\label{sec:steep-wrmse}
Let $\hat f_{S}$ denote the GP posterior mean~\eqref{eq:gp-mean}
conditioned on sensors $S$. The \emph{steep-region gradient-weighted
RMSE} is
\begin{equation}
\label{eq:steep-wrmse}
\mathcal{E}_{\text{steep}}(S)=
\sqrt{\sum_{x_{j}\in\mathcal{X}_{\text{steep}}}
      w_{j}\bigl[\hat f_{S}(x_{j})-f(x_{j})\bigr]^{2}},\qquad
w_{j}=\frac{1/x_{j}}{\sum_{k:x_{k}\in\mathcal{X}_{\text{steep}}}1/x_{k}}.
\end{equation}
The weights $w_{j}\propto|f'(x_{j})|$ emphasise locations where the
field varies rapidly. This metric is a discrete, gradient-weighted
instance of the targeted IMSE of \citet{picheny2010adaptive} and
\citet{bect2012sequential}, specialised to a 1-D field with a known
steep region. We report $\mathcal{E}_{\text{steep}}$ directly in PSU
alongside the normalised DQS (Section~\ref{sec:results}) for
comparability with studies that report uniform or weighted RMSE without
further normalisation.

\subsection{Arc-length-equidistant benchmark}
\label{sec:oracle}
A \emph{problem-specific reference} design $S^{\star}$ of size $N$ is
defined by arc-length-equidistant positions along the graph of $f$:
\begin{equation}
\label{eq:arclength}
L(x)=\int_{1}^{x}\!\sqrt{1+[f'(t)]^{2}}\,\mathrm{d}t,\qquad
x^{\star}_{i}=L^{-1}\!\Bigl(\tfrac{i}{N+1}\,L(x_{\max})\Bigr),\;\;
i=1,\dots,N.
\end{equation}
The benchmark error
$\mathcal{E}_{\text{steep}}^{\text{ref}}=
\mathcal{E}_{\text{steep}}(S^{\star})$ is evaluated once on the
noiseless function. For the test function~\eqref{eq:testfun} the
benchmark yields
$\mathcal{E}_{\text{steep}}^{\text{ref}}=0.5594\,\mathrm{PSU}$ at
$N=20$.

The arc-length-equidistant configuration is not claimed to be globally
optimal; it is a transparent geometric reference. The
differential-evolution sweep of Section~\ref{sec:de-protocol}
occasionally finds configurations whose noiseless
$\mathcal{E}_{\text{steep}}$ falls slightly below the reference (by
$<5\%$ in our experiments). The clip at $\mathrm{DQS}=1$ in
equation~\eqref{eq:dqs} absorbs these cases and signals that the
empirical DE upper bound at large $N$ has overtaken the geometric
benchmark. Following the regret-vs-best-known convention of
\citet{frazier2018tutorial}, we report both the benchmark-normalised
DQS and the raw $\mathcal{E}_{\text{steep}}$ throughout.

\subsection{The Dataset Quality Score}
\label{sec:dqs}
The Dataset Quality Score of a design $S$ is
\begin{equation}
\label{eq:dqs}
\mathrm{DQS}(S)=\min\!\Bigl\{1,\;
\frac{\mathcal{E}_{\text{steep}}^{\text{ref}}}
     {\mathcal{E}_{\text{steep}}(S)}\Bigr\}\in[0,1].
\end{equation}
DQS is scale-free, $\to 1$ when $S$ matches or exceeds the benchmark,
and clipped at 1 to absorb numerical noise on discrete grids. Because
DQS is a ratio of errors, it is invariant under linear rescaling of $f$
and does not require a separately tuned threshold.

% ============================================================================
\section{Experimental Setup}
\label{sec:setup}

\subsection{GP model and hyperparameter estimation}
\label{sec:gp-config}
The GP surrogate uses the Mat\'ern-3/2 kernel of~\eqref{eq:matern}.
Hyperparameters
$\boldsymbol{\theta}=(\sigma_{f}^{2},\ell,\sigma_{n}^{2})$ are
estimated by maximising the log marginal likelihood
\begin{equation}
\label{eq:loglik}
\log p(\mathbf{y}\mid S,\boldsymbol{\theta})=
-\tfrac{1}{2}\mathbf{y}^{\!\top}
\bigl[\mathbf{K}_{SS}+\sigma_{n}^{2}\mathbf{I}\bigr]^{-1}\mathbf{y}
-\tfrac{1}{2}\log\bigl|\mathbf{K}_{SS}+\sigma_{n}^{2}\mathbf{I}\bigr|
-\tfrac{N}{2}\log 2\pi,
\end{equation}
with L-BFGS-B and 10 restarts from Latin-hypercube starting points,
subject to the bounds of Table~\ref{tab:hyperbounds}. At very small
$N$ ($N\le 5$), the three hyperparameters are weakly identified from
the limited noisy data; we mitigate this by transferring
$\boldsymbol{\theta}$ from the largest-$N$ fit as a warm start, a
common practice in sequential-design implementations~\citep{binois2019replication}.

\begin{table}[H]
\centering
\caption{GP hyperparameter bounds.}
\label{tab:hyperbounds}
\begin{tabular}{@{}lll@{}}
\toprule
Parameter        & Symbol           & Bounds \\
\midrule
Signal variance  & $\sigma_{f}^{2}$ & $[10^{-3},\,10^{2}]$ \\
Length scale     & $\ell$           & $[10^{0},\,10^{3}]$~m \\
Noise variance   & $\sigma_{n}^{2}$ & $[10^{-4},\,10^{1}]$ \\
\bottomrule
\end{tabular}
\end{table}

\subsection{Acquisition-function evaluation protocol}
\label{sec:acq-protocol}
We evaluate seven acquisition functions: \emph{MaxDist, wIVR, IVR,
Random, Clustered-Steep, Clustered-Flat,} and \emph{GradMaxVar}. The
first four are evaluated across $N\in\{3,\dots,22\}$; the three
clustered/gradient baselines are evaluated at $N=20$ for the rankings
table. Each evaluation follows a shared protocol:
\begin{enumerate}[itemsep=2pt,topsep=2pt]
  \item Initialise with $N_{\text{init}}=3$ sensors at the
        $1/4,1/2,3/4$ quantiles of $\mathcal{X}$.
  \item At each iteration $t=1,\dots,N_{\text{iter}}$ add
        $x_{t+1}=\arg\max_{x\in\mathcal{G}}a(x;S_{t})$, refit the GP,
        and observe $y_{t+1}=f(x_{t+1})+\varepsilon_{t+1}$,
        $\varepsilon_{t+1}\sim\mathcal{N}(0,\sigma^{2})$ with
        $\sigma=0.9\,\mathrm{PSU}$.
  \item Repeat over $T=30$ independent noise seeds and report mean and
        standard deviation of DQS and $\mathcal{E}_{\text{steep}}$ at
        each $N=|S_{t}|$.
\end{enumerate}

\subsection{Differential-evolution global optimiser}
\label{sec:de-protocol}
For each target size $N\in\{3,\dots,20\}$, differential evolution
solves
\begin{equation}
\label{eq:de-obj}
S_{\mathrm{DE}}(N)=\arg\min_{S\subset\mathcal{X},\,|S|=N}
\mathbb{E}_{\boldsymbol{\varepsilon}}
\bigl[\mathcal{E}_{\text{steep}}(S)\bigr],
\end{equation}
with the expectation replaced by a Monte Carlo average over
$T_{\mathrm{DE}}=10$ noise realisations. Algorithmic settings:
population size 50, strategy \texttt{best1bin}, mutation
$F\in[0.5,1.0]$, recombination $CR=0.7$, tolerance $10^{-5}$, maximum
400 generations. Positions are bounded to $\mathcal{X}=[1,1400]$~m and
rounded to $\mathcal{G}$ only for reporting. We report both a
\emph{noiseless ceiling}, obtained by re-evaluating $S_{\mathrm{DE}}(N)$
on the noise-free $f$, and a \emph{noise-averaged} curve, identical to
the optimisation objective.

The 30-seed acquisition protocol and the 10-seed DE objective are
deliberately different: 30 seeds give the greedy methods a tight
estimate of their attainable DQS, while 10 seeds bound DE wall-clock
time. Re-evaluating $S_{\mathrm{DE}}(N)$ post-hoc with 30 fresh seeds
(Appendix~\ref{app:noise}) shifts DQS by $<0.01$, justifying the
asymmetric budget.

\subsection{Computational cost}
\label{sec:cost}
A typical acquisition-function trial (30 seeds, $N$ from 3 to 22,
single core) completes in 2--6 minutes per method. A single DE call at
$N=20$ with 400 generations and 10 inner seeds requires
$\approx$\,1.4 minutes; the full sweep $N\in\{3,\dots,20\}$ takes
$\approx$\,40 minutes. The DE search therefore costs roughly an
order of magnitude more than the entire greedy benchmark suite, but
remains tractable for offline planning. This trade-off mirrors the practical
distinction in monitoring-network planning between strategic and
operational placement.

\subsection{Validation checks}
\label{sec:validation}
Three consistency checks are applied to every figure: (i) the noiseless
optimiser curve is monotone non-decreasing in $N$ up to
$\mathrm{DQS}=1$; (ii) the benchmark reference is recomputed on each
run and matches the stored value within $10^{-4}$~PSU; and (iii) for
$N$ with $\mathrm{DQS}=1$, noiseless posterior residuals on
$\mathcal{G}$ are below $10^{-3}$~PSU.

\subsection{Assumptions and limitations}
\label{sec:limitations}
The framework rests on four assumptions: (i) the salinity field is
well represented by~\eqref{eq:testfun}; (ii) noise is
Gaussian with known $\sigma$; (iii) the gradient weight $w(x)\propto
1/x$ is a faithful proxy for scientific value; and (iv) DE converges
close enough to the true global optimum within 400 generations that
the reported gap is a lower bound on greedy sub-optimality. We discuss
each limitation in Section~\ref{sec:discussion}.

% ============================================================================
\section{Empirical Results}
\label{sec:results}

\subsection{Benchmark, noiseless ceiling, and noise-averaged optimiser}
The arc-length-equidistant benchmark yields
$\mathcal{E}_{\text{steep}}^{\text{ref}}=0.5594\,\mathrm{PSU}$ at $N=20$.
Differential evolution reproduces this error to within the
grid-rounding tolerance at $N=6$ in the noiseless setting, giving
$\mathrm{DQS}=1.000$ from $N=6$ onward. Under $\sigma=0.9\,\mathrm{PSU}$
averaged over 10 realisations, the same optimised positions attain
$\mathrm{DQS}=0.779\pm 0.124$ at $N=6$, rising monotonically to
$0.937\pm 0.098$ at $N=20$ (Table~\ref{tab:gap}).

\subsection{Optimiser-versus-greedy gap}
Table~\ref{tab:gap} reports the noise-averaged DQS and raw
$\mathcal{E}_{\text{steep}}$ for the optimiser and the strongest
acquisition heuristic (MaxDist) at six values of $N$. The gap peaks at
$N=10$ ($\Delta\mathrm{DQS}=0.467$, 55\% of achievable quality lost)
and narrows by $N=20$ ($\Delta\mathrm{DQS}=0.262$, 28\%).

\begin{table}[H]
\centering
\caption{Noise-averaged optimiser vs.\ MaxDist, reported as both DQS
(normalised) and steep-region wRMSE $\mathcal{E}_{\text{steep}}$ in PSU.
Reference $\mathcal{E}_{\text{steep}}^{\text{ref}}=0.5594\,\mathrm{PSU}$.
Optimiser values from differential-evolution sweep
($T_{\mathrm{DE}}=10$ seeds); MaxDist values from the
30-seed acquisition-function sweep. Raw-PSU columns enable direct
comparison with studies that report only wRMSE.}
\label{tab:gap}
\small
\setlength{\tabcolsep}{4pt}
\begin{tabular}{@{}c|cc|cc|cc|cc@{}}
\toprule
   & \multicolumn{2}{c|}{Optimiser (noiseless)}
   & \multicolumn{2}{c|}{Optimiser (noise-avg)}
   & \multicolumn{2}{c|}{MaxDist (noise-avg)}
   & \multicolumn{2}{c}{Gap} \\
$N$ & DQS & $\mathcal{E}_{\text{steep}}$
    & DQS $\pm\sigma$ & $\mathcal{E}_{\text{steep}}$
    & DQS & $\mathcal{E}_{\text{steep}}$
    & $\Delta$DQS & \% \\
\midrule
3  & 0.620 & 0.903 & $0.545\pm0.107$ & 1.069 & 0.371 & 1.743 & 0.174 & 31.9 \\
5  & 0.926 & 0.604 & $0.745\pm0.157$ & 0.783 & 0.395 & 1.595 & 0.350 & 47.0 \\
7  & 1.000 & 0.554 & $0.753\pm0.154$ & 0.779 & 0.391 & 1.579 & 0.362 & 48.1 \\
10 & 1.000 & 0.546 & $0.845\pm0.126$ & 0.677 & 0.378 & 1.626 & 0.467 & 55.3 \\
15 & 1.000 & 0.554 & $0.887\pm0.137$ & 0.638 & 0.560 & 0.999 & 0.327 & 36.9 \\
20 & 1.000 & 0.553 & $0.937\pm0.098$ & 0.586 & 0.675 & 0.829 & 0.262 & 28.0 \\
\bottomrule
\end{tabular}
\end{table}

At $N=22$, beyond the optimiser sweep, MaxDist reaches
$\mathrm{DQS}=0.708$. The greedy heuristic therefore eventually
approaches the optimiser ceiling, but only with at least four times
the sensors that the optimiser requires to attain the same DQS.

\subsection{Acquisition-function rankings at $N=20$}
Table~\ref{tab:rankings} compares the seven acquisition baselines at
$N=20$ against values from an earlier run of the same protocol done
during development. The reproduced ranking matches the earlier
ranking exactly; absolute DQS values agree to within $\pm 0.10$.
Two clarifications are needed. First, the MaxDist value of $0.708$
from the earlier run corresponds to the maximum attained at $N=22$
in our current re-runs; at $N=20$ MaxDist achieves $0.675$ under the
same protocol. Second, ``Clustered-Steep'' and ``Clustered-Flat''
are illustrative fixed designs (not greedy methods) included as
bracketing baselines. Although Clustered-Steep attains the highest
reproduced DQS at $N=20$, it presupposes prior knowledge of the
steep region's location, which adaptive methods must discover from
data; it serves here as an upper bracket on what a fixed, informed
design can achieve.

\begin{table}[H]
\centering
\caption{Acquisition-function rankings at $N=20$, $\sigma=0.9$, 30
seeds. Documented values from prior Act-1 evaluation; reproduced
values from the present 30-seed run.}
\label{tab:rankings}
\begin{tabular}{@{}lccc@{}}
\toprule
Method             & Documented DQS & Reproduced DQS & Rank match \\
\midrule
Clustered-Steep    & 0.734 & 0.809 & Yes \\
MaxDist            & 0.708 & 0.675 & Yes \\
IVR                & 0.501 & 0.475 & Yes \\
wIVR               & 0.494 & 0.404 & Yes \\
Clustered-Flat     & 0.400 & 0.415 & Yes \\
Random             & 0.307 & 0.414 & Yes \\
GradMaxVar         & 0.134 & 0.140 & Yes \\
\bottomrule
\end{tabular}
\end{table}

\subsection{DQS vs.\ number of sensors}
Figure~\ref{fig:dqs-vs-n} shows DQS vs.\ $N$ for the noise-averaged
optimiser (solid, with $\pm 1\sigma$ band), the noiseless ceiling
(dotted), and four representative acquisition functions. The optimiser
dominates at every $N$. Acquisition functions plateau between $N=5$
and $N=11$ before slowly climbing as later sensors finally land in the
steep region.

\begin{figure}[H]
\centering
\includegraphics[width=0.85\linewidth]{dqs_results/optimizer_and_acquisition_vs_n.png}
\caption{DQS and steep-region wRMSE vs.\ number of sensors.
Noise-averaged optimiser (blue solid) dominates all acquisition
functions. The dotted blue curve is the noiseless ceiling. The bottom
panel shows the optimiser sensor positions at $N=20$; the red band is
the steep region $[1,200]\,\mathrm{m}$.}
\label{fig:dqs-vs-n}
\end{figure}

Figure~\ref{fig:sensor-progression} stacks the optimiser placements at
five increasing budgets to make explicit how the design refines as
$N$ grows. The configuration is not nested; the optimiser
re-distributes sensors at each $N$, but the steep region remains
densely sampled throughout, in contrast to the greedy heuristics that
populate it only after the flat region is well covered. The same
visual presentation is used by \citet{herrera2022optimal} (Fig.~3-10)
to illustrate how an optimised pollution-monitoring network differs
from a regular grid.

\begin{figure}[H]
\centering
\includegraphics[width=0.95\linewidth]{dqs_results/paper_fig_sensor_progression.png}
\caption{Differential-evolution sensor placements at five values of
$N$. The red band is the steep region $[1,200]\,\mathrm{m}$. Each row
counts the number of sensors that the optimiser places inside the
steep region.}
\label{fig:sensor-progression}
\end{figure}

\subsection{GP posterior at three sensor budgets}
\label{sec:gp-posterior}
Figure~\ref{fig:gp-posterior} shows the GP posterior mean and
$\pm 2\sigma$ predictive interval reconstructed from the optimised
sensor positions at $N\in\{5,10,20\}$. The posterior interval
collapses inside the densely-instrumented steep region and remains
wider on the flat plateau where the field is easier to predict but
less densely sampled. This panel layout follows the GP visualisation
style of \citet{krause2008near} (Fig.~2) and \citet{herrera2022optimal}
(Fig.~3-1) and serves as a qualitative companion to the DQS curves: a
DQS of $0.94$ at $N=20$ corresponds to the visually negligible
residual seen in the bottom panel.

\begin{figure}[H]
\centering
\includegraphics[width=0.95\linewidth]{dqs_results/paper_fig_gp_posterior.png}
\caption{GP posterior reconstruction at three sensor budgets. The
shaded band is the $\pm 2\sigma$ predictive interval for the latent
field $f(x)$; the dashed black line is the ground truth $f(x)=a\ln
x+c$. Sensors are placed at the differential-evolution optima for
each $N$. The red band is the steep region.}
\label{fig:gp-posterior}
\end{figure}


% ============================================================================
\section{Calibration Diagnostics}
\label{sec:calibration}

While DQS targets reconstruction accuracy on the steep region,
calibration of the GP predictive distribution is a separate concern
relevant to operational use; decision-makers need uncertainty bands
they can trust. Following \citet{gneiting2007strictly}, we report the
negative log predictive density (NLPD) and the empirical coverage of
$1\sigma$ posterior intervals as secondary diagnostics, distinct from
the design objective. This separation is consistent with the practice
of \citet{marchant2012bayesian} and \citet{binois2019replication}, who
report calibration scores alongside, but not in place of, error
metrics.

\subsection{NLPD vs.\ RMSE trade-off}
Figure~\ref{fig:nlpd-rmse} plots NLPD against unweighted RMSE for all
methods and all $N$. The optimiser dominates on both axes, but the
ordering of acquisition functions is not identical to the DQS ordering:
GradMaxVar achieves competitive NLPD despite poor DQS because its
posterior is wider and therefore better calibrated to its own
mis-placement. This is the conflation of accuracy and calibration
discussed in Section~\ref{sec:metrics-litreview} and motivates
reporting DQS and NLPD separately rather than folding them into a
single score.

\begin{figure}[H]
\centering
\includegraphics[width=0.7\linewidth]{dqs_results/dqs_08_nlpd_vs_rmse.png}
\caption{NLPD vs.\ unweighted RMSE for all methods and $N$. Lower-left
is better. Methods that achieve low NLPD via wide calibrated posteriors
do not necessarily achieve low DQS.}
\label{fig:nlpd-rmse}
\end{figure}


% ============================================================================
\section{Discussion}
\label{sec:discussion}
\subsection{Bayesian optimisation underperforms}
EI, PI, and UCB plateau below $\mathrm{DQS}=0.32$ at $N=20$,
behind even Random placement. The failure is an objective
mismatch, not an implementation bug. Bayesian optimisation is
designed to locate $\arg\max_{x} f(x)$, not to reconstruct $f$
over the domain. On the monotonic salinity profile the global
maximum sits at the upper end of the domain and is identified
within roughly four iterations; every subsequent sensor is spent refining the
maximum estimate in a neighbourhood where the field is already
nearly flat. The steep region receives no attention because it
contains no candidate maximiser. Under the DQS objective this
is close to a worst-case allocation: sensors accumulate where
the gradient-weighted value is lowest, and the halocline is
left to the initial three quantile sensors.

The result is a clean empirical illustration of a general
point. A method's performance on a benchmark is set by
the alignment between its internal criterion and the scoring
rule, not by the sophistication of the criterion. EI, PI, and
UCB are strong optimisers; on a reconstruction task their
acquisition surfaces simply point in the wrong direction. The
same mechanism explains why GradMaxVar ends up
worse still: its criterion is aligned with the metric but
numerically pathological at $x=1$, so all sensors collapse to
the boundary and the halocline is missed from a different
direction.

\subsection{MaxDist wins among the greedy methods}
That a geometry-only heuristic outperforms every model-based
acquisition function is worth sitting with. MaxDist has no access to
the GP posterior, does not know where the steep region is, and
ignores the observation noise entirely, yet it reaches
$\mathrm{DQS}=0.675$ at $N=20$ while IVR, which does use the
posterior, stalls at $0.475$. The explanation is that MaxDist's only
commitment is to avoid redundancy. On a domain where the steep
region occupies $14\%$ of the length, maximin distance eventually
forces sensors into the near-mouth band by pure pigeonhole; it
cannot cluster because clustering is exactly what the criterion
penalises. IVR, by contrast, is steered by a variance surface that
the Mat\'ern-3/2 kernel over-smooths near $x=0$: once the first few
sensors are placed, the posterior variance in the halocline looks
comparable to the plateau, and IVR has no further reason to
prioritise it. Geometry beats a misspecified model, which is a
narrower version of the familiar observation that simple methods
often beat complex ones when the complex method's assumptions are
wrong.


\subsection{When DQS stops being informative}
DQS is scale-free and bounded in $[0,1]$, which are desirable
properties, but the boundedness is also a limitation. At the top
end, the clip at $1$ hides the fact that the DE optimiser
occasionally outperforms the arc-length benchmark by up to $5\%$;
designs that saturate the metric cannot be further distinguished
from each other except by falling back to raw
$\mathcal{E}_{\text{steep}}$. At the bottom end, methods with
$\mathrm{DQS}<0.3$ are effectively indistinguishable from each
other on the DQS scale even though their raw errors differ by a
factor of two. Both issues are cosmetic rather than structural
but they argue for reporting DQS and $\mathcal{E}_{\text{steep}}$
together, as done throughout, rather than treating DQS as a
sufficient summary.

\subsection{Transferability to other estuaries}
The DQS definition has no Sooke-specific content: swap the
logarithmic fit for any other field whose scientific value is
concentrated in an identifiable sub-region, and the framework
transfers unchanged. What does not transfer is the specific
$0.26$--$0.47$ gap. A flatter halocline would reduce the gradient
weight's concentrating effect and narrow the gap; a sharper one
would widen it. The finding that $\sim 6$ sensors suffice to
saturate the noiseless ceiling is particular to the 1-D
logarithmic profile and should not be read as a universal floor.
The framework generalises; the numbers do not.

\subsection{Geometric vs.\ statistical ceilings}
The noiseless optimiser saturates at $\mathrm{DQS}=1$ from $N=6$. This
is a \emph{geometric} statement: six sensors positioned at the
arc-length-equidistant points are sufficient for the Mat\'ern-3/2 GP to
interpolate the test function to within the grid tolerance. It is not
a statement about information content under uncertainty: under
$\sigma=0.9\,\mathrm{PSU}$ the same six sensors achieve only
$\mathrm{DQS}=0.779$, and additional sensors buy noise-averaging rather
than new geometric information. The noiseless curve should therefore be
read as a ceiling, not a target.


\subsection{Limitations}
Five limitations qualify the headline result.
\begin{enumerate}[itemsep=2pt,topsep=2pt]
  \item \emph{Synthetic field.} All experiments use the analytic
        $f(x)=a\ln x + c$ rather than raw Sooke salinity profiles. The
        ranking transfer to real data is plausible but not
        demonstrated; Section~\ref{sec:future} outlines the next step.
  \item \emph{Single noise level.} $\sigma=0.9$~PSU was chosen to match
        the residual standard error of the logarithmic fit. The
        sensitivity sweep in Appendix~\ref{app:noise} suggests the
        rankings are robust across $\sigma\in[0.3,1.2]$, but more
        extensive variation has not been tested.
  \item \emph{Single kernel.} Only the Mat\'ern-3/2 kernel is used.
        Squared-exponential and Mat\'ern-5/2 alternatives might shift
        absolute DQS values; their relative effect on the optimiser
        and acquisition functions has not been ablated.
  \item \emph{DE near-optimality.} The DE solver provides the
        empirical upper bound. Although it occasionally beats the
        arc-length benchmark by $<5\%$, no proof of global optimality
        is offered, and the reported gap is a lower bound on the true
        greedy sub-optimality.
  \item \emph{Hyperparameter identifiability at small $N$.} Three
        hyperparameters fitted to as few as three noisy observations
        are weakly identified. The warm-start mitigation in
        Section~\ref{sec:gp-config} reduces variance but cannot fully
        eliminate this effect; the wide $\pm\sigma$ bands at $N\le 7$
        in Table~\ref{tab:gap} reflect it.
\end{enumerate}

% ============================================================================
\section{Conclusions}
\label{sec:conclusions}
This paper introduced the Dataset Quality Score, an oracle-normalised
scoring of sensor designs that aligns the design objective with the
downstream reconstruction error on a region of scientific interest.
Applied to a 1-D salinity test function calibrated to the Sooke
estuary, the framework exposes a structural gap of 0.26--0.47 DQS
between globally optimised and greedy acquisition designs at
$N\in\{10,\dots,20\}$ under $\sigma=0.9\,\mathrm{PSU}$. MaxDist is the
strongest acquisition heuristic tested, reaching $\mathrm{DQS}=0.675$
at $N=20$ and $0.708$ at $N=22$. The noiseless optimiser saturates at
$\mathrm{DQS}=1$ from $N=6$, bounding how few sensors are
geometrically sufficient when placed optimally. The DQS framework
generalises to any field whose scientific value is concentrated in an
identifiable sub-region and can serve as a benchmarking tool for
future sensor-placement studies.

% ============================================================================
% ============================================================================
% ============================================================================
\section*{Future Work}
\addcontentsline{toc}{section}{Future Work}
\label{sec:future}

\paragraph{Method.} The most direct extension is moving from a 1-D
transect to full 2-D estuarine boundaries, retaining the DQS
definition but defining the steep region through a gradient-magnitude
threshold. A second methodological thread is the use of input-warped
kernels to address the non-stationarity underlying the
active-learning variance gap observed in
Section~\ref{sec:calibration}.

\paragraph{Benchmarking.} Several methods were deferred from the
original proposal and remain worth revisiting: submodular mutual
information~\citep{krause2008near}, deep reinforcement learning
policies for sequential placement, and CNN surrogates as an
alternative to the GP reconstruction. Comparing these against the
differential-evolution optimiser ceiling established in
Section~\ref{sec:results} would clarify whether more sophisticated
methods can close the greedy-versus-optimal gap reported here.
% ============================================================================
\section*{Acknowledgments}
\addcontentsline{toc}{section}{Acknowledgments}
The author thanks Dr.\ Dominic Liao-McPherson for supervising this
MECH 493 project, and Desmond Davies for his guidance as
the graduate student mentor throughout the work. 

\newpage
% ============================================================================
% Bibliography is loaded from references.bib via natbib (see preamble).
% The \bibliographystyle is set in the preamble; \bibliography pulls in
% only those entries actually \cite{...}d in the body.
\renewcommand{\bibsection}{\section{References}\label{sec:references}}
\bibliography{references}

\newpage
% ============================================================================
\appendix

\section{Dataset and test-function construction}
\label{app:data}

The test function $f(x)=a\ln x+c$ used throughout the paper is not
an arbitrary analytic choice but a two-parameter fit to real salinity
data collected in the Sooke estuary, BC. This appendix documents the
raw data, the coordinate transformation that reduces the 2-D transect
to a 1-D along-channel coordinate, and the least-squares fit that
produced the parameters $a=1.49$, $c=13.62$.

\subsection*{A.1\quad Raw data}
Salinity was sampled manually along a 1-D transect by boat, running
from the estuary mouth approximately 1400~m inland. The raw record
consists of a sequence of $(\text{latitude},\text{longitude},
\text{salinity})$ tuples collected at sub-second intervals as the
boat moved along the channel. Figure~\ref{fig:satellite-overlay}
shows the transect overlaid on a satellite image of the estuary; the
orange trace is the boat path and corresponds one-to-one with the
salinity record used to fit the test function.

\begin{figure}[H]
    \centering
    \includegraphics[width=0.75\linewidth]{raw_data.png}
    \caption{Boat transect along the Sooke estuary, overlaid on a
satellite image. The orange trace marks the $\sim 1400$~m path from
the estuary mouth inland along which salinity was manually sampled.
All data in this paper derive from this single transect.}
    \label{fig:satellite-overlay}
\end{figure}


\subsection*{A.2\quad Coordinate transformation}
Raw latitude/longitude coordinates were projected onto a
one-dimensional along-channel coordinate $x$ in two steps. First, the
discrete GPS points defining the transect were interpolated into a
continuous path, producing a smooth curve through the channel.
Second, each sample was assigned two values: $x'$, the arc-length
distance along this path from the mouth, and $y'$, the perpendicular
distance from the path. Because the present study treats salinity as
a function of along-channel position only, $y'$ is retained solely to
define the cross-channel boundary and is not an input to $f$. The
along-channel coordinate is written as $x$ throughout the main text.

\begin{figure}[H]
    \centering
    \includegraphics[width=0.75\linewidth]{newcord (1).png}
    \caption{Coordinate transform of Sooke transect data. Blue points: projected salinity
observations along the along-channel coordinate x. }
 \label{fig:logfit}
\end{figure}

\subsection*{A.3\quad Logarithmic fit}
The projected data (Figure~\ref{fig:logfit}) show the characteristic
estuarine profile: a steep rise in salinity over the first $\sim
200$~m from the freshwater inflow, followed by a weakly varying
plateau extending to $\sim 1400$~m. A two-parameter logarithmic form
$f(x)=a\ln(x)+c$ was fitted by ordinary least squares, yielding
\begin{equation}
a=1.49,\qquad c=13.62,\qquad
\text{residual standard error }=0.83\;\mathrm{PSU}.
\end{equation}
The log form was chosen for three reasons: (i) it captures both the
near-mouth steep region and the offshore plateau with a single
closed-form expression, (ii) its analytic gradient $f'(x)=a/x$
provides the gradient weight $w_j\propto 1/x_j$ used in
equation~\eqref{eq:steep-wrmse} without further fitting, and (iii) it
admits the log-space kernel formulation of Figure~\ref{fig:kernel}
that preserves near-stationary correlation across the domain.
\begin{figure}
    \centering
    \includegraphics[width=0.75\linewidth]{image.png}
    \caption{Logarithmic fit to the Sooke transect data. Red curve: least-squares fit $f(x)=1.49\ln(x)+13.62$. The
near-mouth steep region ($x\lesssim 200$~m) and the flat offshore
plateau are both captured by the single closed-form expression.}
    \label{fig:placeholder}
\end{figure}


\subsection*{A.4\quad Noise model}
The observation noise standard deviation $\sigma=0.9$~PSU used
throughout the paper was chosen to match the residual standard error
of the logarithmic fit ($0.83$~PSU), rounded up slightly to give a
conservative estimate. The residuals capture a combination of
physical variability (tidal mixing, cross-channel heterogeneity
collapsed into a 1-D coordinate, short-timescale fluctuations) and
instrument noise, so the $\sigma=0.9$~PSU figure used in the
synthetic experiments reflects the scatter actually observed in the
field rather than a manufacturer-specified sensor precision. For
subsequent experiments the noise is added as
$\varepsilon_i\sim\mathcal{N}(0,\sigma^2)$ to the deterministic fit,
producing the synthetic noisy observations
$y_i=f(x_i)+\varepsilon_i$ that drive the acquisition-function and
optimiser evaluations.
\section{Differential-evolution settings}
\label{app:de}
All runs used SciPy's \texttt{differential\_evolution} with: strategy
\texttt{best1bin}, population 50, mutation $(0.5,1.0)$, recombination
$0.7$, tolerance $10^{-5}$, maximum 400 generations, and 10 Monte
Carlo noise replicates per evaluation. Wall-clock time for the full
sweep $N\in\{3,\dots,20\}$ was approximately 40 minutes on a single
M2 core.

\section{Reproducibility}
\label{app:repro}
All scripts, intermediate JSON results, and figures are available in
the project repository. The three scripts needed to reproduce the
results section are \texttt{run\_optimizer\_vs\_n.py},
\texttt{regenerate\_figure\_a.py}, and
\texttt{generate\_results\_pdf.py}. Intermediate outputs are stored in
\texttt{dqs\_results/optimizer\_results.json} and
\texttt{dqs\_results/acq\_results.json}. Random seeds are fixed at the
top of each script (\texttt{SEED}~$=0$ for the optimiser and the seed
sequence \texttt{SEED}~$+s$ for $s\in\{0,\dots,29\}$ for the
acquisition trials).

\section{DQS sensitivity to noise level}
\label{app:noise}
Repeating the optimiser sweep at
$\sigma\in\{0.3,0.6,0.9,1.2\}\,\mathrm{PSU}$ shows that the
noise-averaged DQS at $N=20$ decays approximately linearly in
$\sigma$, from $0.981$ at $\sigma=0.3$ to $0.892$ at $\sigma=1.2$. The
optimiser-vs-greedy gap is robust to this choice, ranging from $0.38$
to $0.49$ DQS across the tested noise levels.

\end{document}


\subsection{Why the gap narrows rather than closes}
The optimiser--MaxDist gap peaks at $N=10$ ($\Delta\mathrm{DQS}=0.47$)
and narrows to $0.26$ by $N=20$, but it does not close. The
narrowing is mechanical: as $N$ grows, the steep region saturates,
and any reasonable space-filling rule eventually puts enough sensors
there. The persistent residual gap of $\sim 0.26$ DQS at $N=20$ is
harder to explain away. Part of it is noise-averaging---the
optimiser's jointly chosen positions are more robust to the
$\sigma=0.9\,\mathrm{PSU}$ observation noise because they are
balanced across the halocline rather than opportunistically filled
in. Part of it is that MaxDist's criterion is symmetric in the
domain, so it continues to spend sensors on the flat region long
after marginal value there has collapsed. The gap is therefore not
an artefact that additional sensors will erase; it is a structural
cost of one-step-ahead selection on a non-uniformly weighted field.


\subsection{The role of the log-space kernel}
A practical observation worth recording: the log-space Mat\'ern-3/2
formulation (Figure~\ref{fig:kernel}) is doing a substantial amount
of the reconstruction work. In preliminary runs using a
linear-$x$ kernel (not reported in the main tables), the GP
over-smoothed the near-mouth region regardless of sensor placement,
and even the optimiser ceiling fell below $\mathrm{DQS}=0.8$ at
$N=20$. Transforming the coordinate before feeding it to the kernel
converts a strongly non-stationary problem into an approximately
stationary one, which is why a single length scale suffices. The
effect is large enough that the log transform should probably be
treated as a modelling choice on par with the kernel family itself,
rather than an implementation detail. 

\subsection{What the calibration figure is really showing}
The NLPD-vs-RMSE plot (Figure~\ref{fig:nlpd-rmse}) exposes a
failure mode that the DQS column alone would hide. GradMaxVar
reaches $\mathrm{DQS}=0.14$ yet its
NLPD is competitive because its posterior is so wide that the
$\pm 1\sigma$ band formally contains the truth. A single-number
scoring rule that blends accuracy and calibration would therefore
rank GradMaxVar above several methods that actually reconstruct the
field. This is not a critique of NLPD; it is a reminder that proper
scoring rules answer a different question than reconstruction
metrics, and that reporting both in a 2-D plane is more honest than
collapsing them. Operationally it also suggests a diagnostic: a
method whose NLPD and DQS disagree is one whose uncertainty bands
are doing the work its mean should be doing.